\documentclass[12pt]{article}

\usepackage{amsmath}
\usepackage{amssymb}
\usepackage{geometry}
\usepackage{setspace}

\geometry{letterpaper, margin=1in}
\setstretch{1.15}

\begin{document}

\title{\textbf{Exploring the Identity: \\ $G^* = \sqrt{2\pi} \cdot \theta_3(e^{-\pi})^2$}}
\author{}
\date{}
\maketitle

This is a mathematically beautiful identity that sits at the deep intersection of modular forms, elliptic integrals, continuous probability, and discrete geometry.

By unpacking the expression $G^* = \sqrt{2\pi} \cdot \theta_3(e^{-\pi})^2$ step-by-step, we can reveal its exact closed-form value, its elegant physical interpretations, and its direct relationship to \textbf{Gauss's constant}.

\section*{1. Evaluating the Jacobi Theta Function}
The core of this identity is the Jacobi theta function $\theta_3(q)$, which is defined as an infinite sum of a Gaussian sequence over the integers:
\begin{equation}
    \theta_3(q) = \sum_{n=-\infty}^\infty q^{n^2}
\end{equation}

When we evaluate this at $q = e^{-\pi}$, we are looking at the highly symmetric ``self-dual'' point in the complex plane ($\tau = i$). Because the Gaussian curve is its own Fourier transform, evaluating it at this precise parameter yields a known analytical exact value in terms of the Euler Gamma function ($\Gamma$):
\begin{equation}
    \theta_3(e^{-\pi}) = \frac{\pi^{1/4}}{\Gamma(3/4)}
\end{equation}

To make this easier to work with, we can apply \textbf{Euler's reflection formula}, $\Gamma(z)\Gamma(1-z) = \frac{\pi}{\sin(\pi z)}$. Setting $z = 1/4$:
\begin{align}
    \Gamma(1/4)\Gamma(3/4) &= \frac{\pi}{\sin(\pi/4)} = \pi\sqrt{2} \\
    \implies \Gamma(3/4) &= \frac{\pi\sqrt{2}}{\Gamma(1/4)}
\end{align}

Substituting this back into our evaluated theta function gives:
\begin{equation}
    \theta_3(e^{-\pi}) = \frac{\pi^{1/4} \Gamma(1/4)}{\pi\sqrt{2}} = \frac{\Gamma(1/4)}{\sqrt{2}\pi^{3/4}}
\end{equation}

\section*{2. Squaring and Multiplying by $\sqrt{2\pi}$}
Next, the identity squares this theta value. Geometrically, squaring the 1D theta function gives us the sum of a Gaussian over a 2D square integer lattice ($\mathbb{Z}^2$):
\begin{equation}
    \theta_3(e^{-\pi})^2 = \sum_{m=-\infty}^\infty \sum_{n=-\infty}^\infty e^{-\pi(m^2+n^2)} = \frac{\Gamma(1/4)^2}{2\pi^{3/2}}
\end{equation}

Finally, we multiply by $\sqrt{2\pi}$ as defined in the identity to uncover the exact closed algebraic form of $G^*$:
\begin{align}
    G^* &= \sqrt{2\pi} \cdot \frac{\Gamma(1/4)^2}{2\pi^{3/2}} \\
        &= \frac{\sqrt{2}\pi^{1/2} \Gamma(1/4)^2}{2\pi^{3/2}}
\end{align}

Simplifying the constants ($\pi^{1/2} / \pi^{3/2} = 1/\pi$) yields the true nature of the expression:
\begin{equation}
    \boxed{ G^* = \frac{\Gamma(1/4)^2}{\pi\sqrt{2}} \approx 2.9586751\dots }
\end{equation}

\section*{3. The Bridge to Famous Historical Constants}
Because $G^*$ is driven by $\Gamma(1/4)^2$, it acts as a mathematical ``hub'' that algebraically connects several legendary 19th-century math constants:

\begin{itemize}
    \item \textbf{Gauss's Constant ($G$):} In 1799, Carl Friedrich Gauss discovered that the reciprocal of the arithmetic-geometric mean (AGM) of $1$ and $\sqrt{2}$ yields a fundamental constant:
    \begin{equation*}
        G = \frac{1}{\operatorname{agm}(1, \sqrt{2})} = \frac{\Gamma(1/4)^2}{2\sqrt{2}\pi^{3/2}} \approx 0.8346
    \end{equation*}
    If we extract the Gamma values from our closed form, the identity reveals itself as a scaled version of Gauss's constant:
    \begin{equation}
        \boxed{ G^* = 2\sqrt{\pi} \cdot G }
    \end{equation}

    \item \textbf{The Lemniscate Constant ($\varpi$):} Just as $\pi$ governs the perimeter of a circle, $\varpi \approx 2.622$ governs the perimeter of the Bernoulli Lemniscate (the $\infty$-shaped curve). Because $\varpi = \pi G$, the constant scales directly with the arc length of the Lemniscate:
    \begin{equation}
        \boxed{ G^* = \frac{2}{\sqrt{\pi}} \cdot \varpi }
    \end{equation}

    \item \textbf{The Beta Function:} The Beta function evaluated at one-quarter is $B(1/4, 1/4) = \frac{\Gamma(1/4)^2}{\sqrt{\pi}}$. This allows us to write the identity elegantly normalized by the square root of $2\pi$:
    \begin{equation}
        \boxed{ G^* = \frac{1}{\sqrt{2\pi}} B(1/4, 1/4) }
    \end{equation}

    \item \textbf{Complete Elliptic Integrals:} The complete elliptic integral of the first kind evaluated at $k=1/\sqrt{2}$ yields $K(1/\sqrt{2}) = \frac{\Gamma(1/4)^2}{4\sqrt{\pi}}$. Thus:
    \begin{equation}
        \boxed{ G^* = \frac{2\sqrt{2}}{\sqrt{\pi}} K(1/\sqrt{2}) }
    \end{equation}
\end{itemize}

\section*{4. The Continuous $\times$ Discrete Interpretation}
Perhaps the most poetic way to look at this identity is as a direct collision of continuous probability and discrete lattice mathematics.

Look at the two distinct parts of the original formula:
\begin{enumerate}
    \item $\sqrt{2\pi}$ is the exact total area/normalization factor of the \textbf{1D continuous standard Gaussian integral}: $\int_{-\infty}^\infty e^{-x^2/2} \, dx$.
    \item $\theta_3(e^{-\pi})^2$ is the \textbf{2D discrete Gaussian partition function} acting over a square integer grid.
\end{enumerate}

Therefore, the identity can conceptually be rewritten as:
\begin{equation}
    G^* = \left( \int_{-\infty}^\infty e^{-x^2/2} \, dx \right) \times \left( \sum_{(m,n) \in \mathbb{Z}^2} e^{-\pi(m^2+n^2)} \right)
\end{equation}

\vspace{0.5cm}
\noindent \textbf{In Summary:} \\
The identity miraculously takes the continuous volume of a normal probability distribution ($\sqrt{2\pi}$), multiplies it by the discrete heat kernel of a 2D integer lattice ($\theta_3(e^{-\pi})^2$), and neatly produces a fundamental constant $\left(\frac{\Gamma(1/4)^2}{\pi\sqrt{2}}\right)$ that is permanently tethered to Gauss's arithmetic-geometric mean.

\end{document}
